\documentclass{resume} % Use the custom resume.cls style

\usepackage[left=0.4in,top=0.4in,right=0.4in,bottom=0.4in]{geometry}
\usepackage[utf8]{inputenc}
\usepackage[T2A]{fontenc}
\usepackage[russian]{babel}
\usepackage{fontawesome5}
\usepackage{graphicx}

\name{Караганов Павел}
\address{Java Backend Разработчик (Стажёр)}
\address{
    \href{mailto:466113@niuitmo.ru}{\faEnvelope\hspace{1}466113@niuitmo.ru} \hspace{2ex}
    \href{https://t.me/Antichnoste}{\faTelegram\hspace{1}@Antichnoste} \hspace{2ex}
    \href{https://github.com/Antichnoste}{\faGithub\hspace{1}github.com/Antichnoste}
}

\begin{document}

\begin{rSection}{О себе}
Студент 3 курса ИТМО (поступил по БВИ). Увлекаюсь промышленной backend-разработкой на Java и проектированием отказоустойчивых систем. Совмещаю сильную алгоритмическую базу (олимпиады по спортивному программированию) с практическим опытом работы с современным стеком: Spring Boot, PostgreSQL, Docker, системы мониторинга. Быстро обучаюсь, пишу эффективный код и стремлюсь попасть в Т-Академию для развития навыков построения надежных высоконагруженных (Highload) сервисов.
\end{rSection}

%----------------------------------------------------------------------------------------
%	EDUCATION SECTION
%----------------------------------------------------------------------------------------
\begin{rSection}{Образование}
{\bf Программная Инженерия}, бакалавриат ИТМО (Факультет ПИиКТ) \hfill {2024 - 2028}
\begin{itemize}
    \itemsep -3pt {} 
    \item Поступил без вступительных испытаний (БВИ) по результатам перечневых олимпиад.
    \item Профильные дисциплины: Алгоритмы и структуры данных, Математический анализ, Базы данных.
\end{itemize}
\end{rSection}

%----------------------------------------------------------------------------------------
% TECHINICAL STRENGTHS	
%----------------------------------------------------------------------------------------
\begin{rSection}{Навыки}
\begin{tabular}{ @{} >{\bfseries}l @{\hspace{4ex}} l }
Языки и Основы: & Java (Core, Collections, Multithreading, Stream API), SQL, C++ (алгоритмы) \\
Фреймворки и БД: & Spring (Boot, Security/JWT, Data), PostgreSQL, Liquibase \\
Архитектура и Сборка: & REST API, OpenAPI (Swagger), Maven (настройка сборки) \\
DevOps и Мониторинг: & Docker, Docker Compose, GitHub Actions (CI/CD), Prometheus, Grafana \\
Тестирование: & JUnit, Mockito, Yandex Tank (нагрузочное тестирование)
\end{tabular}
\end{rSection}

%----------------------------------------------------------------------------------------
%	ACHIEVEMENTS SECTION
%----------------------------------------------------------------------------------------
\begin{rSection}{Олимпиады и Достижения}
\begin{itemize}
    \itemsep -3pt {}
    \item \textbf{Спортивное программирование:} Призёр перечневых олимпиад, участник регионального этапа ВсОШ. Опыт оптимизации алгоритмов по времени и памяти.
    \item \textbf{Аналитическое мышление (Шахматы):} 1-й взрослый разряд (рейтинг РШФ 2000). Умение строить гипотезы и находить нестандартные решения в сложных системах.
\end{itemize}
\end{rSection} 

%----------------------------------------------------------------------------------------
%	WORK EXPERIENCE / PROJECTS SECTION
%----------------------------------------------------------------------------------------
\begin{rSection}{Практический опыт и Проекты}

\item \textbf{Разработка клиент-серверного веб-приложения} \hfill \textit{Java, Spring Boot, DevOps}
\begin{itemize}
    \itemsep -3pt {}
    \item Спроектировал REST API, настроил безопасность через \textbf{Spring Security (JWT)} и задокументировал контракты с помощью \textbf{OpenAPI (Swagger)}.
    \item Настроил интеграцию с базой данных \textbf{PostgreSQL} (управление миграциями через \textbf{Liquibase}) и упаковал инфраструктуру в \textbf{Docker Compose}.
    \item Для контроля надежности прикрутил сбор метрик приложения через \textbf{Prometheus}, настроил визуальные дашборды в \textbf{Grafana} и провел нагрузочное тестирование с помощью \textbf{Yandex Tank}.
    \item Внедрил пайплайн CI/CD в GitHub Actions (автоматическая сборка через \textbf{Maven}, тесты, линтеры).
\end{itemize}

\item \textbf{Участие в хакатоне ЛЦТ (Лидеры цифровой трансформации)} \hfill \textit{Backend / DB Разработчик}
\begin{itemize}
    \itemsep -3pt {}
    \item В команде разработал прототип приложения для решения конкурсного кейса в условиях 48-часового дедлайна.
    \item Отвечал за проектирование структур данных, реализацию бизнес-логики и интеграцию с БД.
\end{itemize}

\item \textbf{\href{https://github.com/Antichnoste/OlProga}{Инженерный репозиторий олимпиадных решений (C++)}} \hfill \textit{Алгоритмы и структуры данных}
\begin{itemize}
    \itemsep -3pt {}
    \item Реализовал эффективные алгоритмы для задач с жесткими ограничениями по времени и памяти. 
    \item Проект оформлен по стандартам промышленной разработки: настроена автоматическая сборка через \textbf{CMake} и локальное тестирование на краевых случаях.
\end{itemize}

\end{rSection} 

\end{document}